% Part: first-order-logic
% Chapter: model-theory
% Section: partial-iso

\documentclass[../../../include/open-logic-section]{subfiles}

\begin{document}

\olfileid{mod}{bas}{dlo}
\section{Urutan Linear Rapet Tanpa Titik Ujung}

\begin{defn}
  \emph{Urutan linear rapet tanpa titik ujung}
  iku !!a{structure} $\Struct{M}$
  kanggo !!{language} kang mung ngemot
  siji !!{predicate}~$<$ kanthi
  rong panggonan lan nyukupi
  ukara-ukara iki:
  \begin{enumerate}
  \item $\lforall[x][\lnot x < x]$;
  \item $\lforall[x][\lforall[y][\lforall[z][(x < y \lif (y < z \lif x
    <z ))]]]$;
  \item $\lforall[x][\lforall[y][(x< y \lor \eq[x][y] \lor y < x)]]$;
  \item $\lforall[x][\lexists[y][x < y]]$;
  \item $\lforall[x][\lexists[y][y < x]]$;
  \item $\lforall[x][\lforall[y][(x < y \lif \lexists[z][(x < z \land
        z < y))]]]$.
 \end{enumerate}
\end{defn}

\begin{thm}\ollabel{thm:cantorQ}
  Sembarang rong urutan linear rapet
  tanpa titik ujung kang !!{enumerable}
  padha isomorfis.
\end{thm}

\begin{proof}
  Upamakna $\Struct{M_1}$ lan
  $\Struct{M_2}$ iku urutan linear rapet
  tanpa titik ujung kang !!{enumerable},
  kanthi ${<_1} = \Assign{<}{M_1}$
  lan ${<_2} =
  \Assign{<}{M_2}$, lan upamakna
  $\PIso{I}$ iku himpunan kabeh
  isomorfisme parsial antarane loro
  struktur iku. $\PIso{I}$ ora kosong
  amarga paling ora $\emptyset \in \PIso{I}$.
  Kita nuduhake yen $\PIso{I}$
  nyukupi sipat Maju-Mundur.
  Mula $\Struct{M_1} \iso[p] \Struct{M_2}$,
  lan teorema ndherek saka
  \olref[pis]{thm:p-isom1}.

  Kanggo nuduhake yen $\PIso{I}$
  nyukupi sipat Maju, jupuken
  $p \in \PIso{I}$.
  Yen fungsi parsial iki kosong,
  saben unsur sumber kang dikarepake
  bisa dipasang karo sembarang unsur
  ranah struktur kapindho.
  Kanggo kasus kang ora kosong, upamakna
  $p(a_i) = b_i$ kanggo
  $i = 1$, \dots,~$n$.
  Tanpa ngilangake umumé,
  upamakna $a_1 <_1 a_2 <_1 \cdots <_1
  a_n$. Yen diwenehi
  $a \in \Domain{M_1}$,
  goleka $b \in \Domain{M_2}$
  kaya mangkene:
  Yen unsur kang arep ditambah
  wis ana ing ranah fungsi,
  gunakna pasangane kang wis ana.
  Ing kasus liyane, gunakna
  salah siji saka telung
  pilihan iki:
  \begin{enumerate}
  \item yen $a <_1 a_1$, jupuken
    $b \in \Domain{M_2}$ saéngga
    $b <_2 b_1$;
  \item yen $a_n <_1 a$, jupuken
    $b \in \Domain{M_2}$ saéngga
    $b_n <_2 b$;
 \item yen $a_i <_1 a <_1 a_{i+1}$
   kanggo sawenehing $i$,
   jupuken $b \in \Domain{M_2}$
   saéngga $b_i <_2 b <_2 b_{i+1}$.
  \end{enumerate}
  Cathetan koreksi sumber OLPL-093:
  bukti Maju uga kudu nangani
  fungsi parsial kosong lan
  unsur kang wis ana ing ranahé;
  telung kasus ketaksamaan
  mung nyakup unsur anyar
  nalika ranahé ora kosong.
  Unsur $b$ kanthi sipat
  kang dikarepake mesthi
  bisa ditemokake amarga
  $\Struct{M_2}$ iku urutan
  linear rapet tanpa titik ujung.
  Tetepna $q = p \cup
  \{ \langle a, b \rangle \}$
  saéngga $q \in \PIso{I}$
  dadi ekstensi $p$
  kang dikarepake.
  Iki mbuktekake sipat Maju.
  Sipat Mundur dibuktekake
  kanthi cara padha.
  Mula $\Struct{M_1} \iso[p]
  \Struct{M_2}$; miturut
  \olref[pis]{thm:p-isom1},
  $\Struct{M_1} \iso
  \Struct{M_2}$.
\end{proof}

\begin{prob}
  Rampungna bukti
  \olref[mod][bas][dlo]{thm:cantorQ}
  kanthi mriksa yen
  $\PIso{I}$ nyukupi
  sipat Mundur.
\end{prob}

\begin{rem}
  Upamakna $\Struct{S}$ sembarang
  urutan linear rapet
  tanpa titik ujung kang
  !!{enumerable}.
  Mula (miturut
  \olref{thm:cantorQ})
  $\Struct{S} \iso
  \Struct{Q}$, kanthi
  $\Struct{Q} = (\Rat, <)$
  iku urutan linear rapet
  tanpa titik ujung kang
  !!{enumerable} lan
  ranahé himpunan $\Rat$
  saka wilangan rasional.
  Saiki gatekna maneh
  !!{structure}~$\Struct{R} =
  (\Real, <)$ saka
  \olref[thm]{remark:R}.
  Kita wis weruh yen ana
  !!a{enumerable}
  !!{structure}~$\Struct{S}$
  saéngga $\Struct{R}
  \elemequiv \Struct{S}$.
  Nanging $\Struct{S}$
  iku urutan linear rapet
  tanpa titik ujung kang
  !!a{enumerable}, mula
  isomorfis (lan mulané
  ekuivalen elementer)
  karo !!{structure}~$\Struct{Q}$.
  Amarga ekuivalensi elementer
  transitif, $\Struct{R}
  \elemequiv \Struct{Q}$.
  (Iki uga bisa dibuktekake
  langsung kanthi nuduhake
  $\Struct{R} \iso[p]
  \Struct{Q}$ nganggo
  argumen Maju-Mundur
  kang padha.)
\end{rem}
\end{document}
